\begin{abstract}
\EmoNet{}은 한국어 감정 응답 생성을 입력 자극 인코딩, 감정 동역학, branch/trajectory 요약, 잠재표현, 스타일 회귀, LLM 표면화로 분해하는 프레임워크다. 이 시스템의 목적은 단순히 응답 품질을 높이는 것보다, 감정 제어 실패가 어느 단계에서 발생하는지 진단 가능한 중간표현을 만드는 데 있다. 최신 cycle에서는 calibrated reference configuration을 도입해 dominant branch 평균 길이를 1.0539에서 70.4684로 늘리고 길이 1 비율을 0.9734에서 0.0948로 낮췄다. 또한 raw trajectory를 직접 읽는 heuristic readout 대신 GPT-5.4 기반 episode interpretation 경로를 추가해 trajectory 수준 affect episode를 더 풍부하게 기술할 수 있게 했다. 그러나 end-to-end generation에서는 여전히 \texttt{stim\_only}가 mean total 3.5404로 가장 높고, \texttt{episode\_trace}는 3.4673, \texttt{emonet\_full}은 3.2478에 머문다. predictor 비교에서도 learned $z \rightarrow s$ 경로는 text baseline보다 강하지 않다. 따라서 현재 \EmoNet{}의 가장 강한 기여는 최종 성능 우위 자체보다, branch collapse 완화와 trajectory interpretation을 통해 감정 처리의 내부 병목을 측정 가능한 형태로 드러냈다는 데 있다.
\end{abstract}
